\section{Results}
\label{sec:results}

\subsection{The refinement paradox: retrieval quality and faithfulness diverge}

Table~\ref{tab:hcpc_main} reports the central comparison from Phase~6. The
results directly test the paper's core hypothesis: that retrieval quality
improvements do not reliably transfer to faithfulness improvements.

\begin{table}[!htbp]
\centering
\caption{HCPC head-to-head comparison (Mistral-7B, chunk = 1024, $k = 3$,
$n = 30$ per condition). The paradox is immediately visible: HCPC-v1
\emph{maintains or improves} retrieval similarity while \emph{degrading}
faithfulness and introducing hallucination.}
\label{tab:hcpc_main}
\small
\begin{tabular}{llccccc}
\toprule
Dataset & Condition & Faith. & Halluc. & Mean Sim & Refine\% & CCS \\
\midrule
SQuAD    & Baseline & 0.7995 & 0.0\%  & 0.634 & ---    & --- \\
SQuAD    & HCPC-v1  & 0.6407 & 10.0\% & 0.624 & 100.0\% & --- \\
SQuAD    & HCPC-v2  & 0.7874 & 0.0\%  & 0.634 & 16.7\% & 1.00 \\
\midrule
PubMedQA & Baseline & 0.5783 & 20.0\% & 0.653 & ---    & --- \\
PubMedQA & HCPC-v1  & 0.5808 & 26.7\% & 0.678 & 90.0\% & --- \\
PubMedQA & HCPC-v2  & 0.5893 & 20.0\% & 0.653 & 73.3\% & 0.57 \\
\bottomrule
\end{tabular}
\end{table}

The paradox manifests differently across datasets, revealing its dependence
on upstream retrieval quality. On SQuAD, where the encoder-corpus alignment
is strong, HCPC-v1 turns an already coherent retrieval set into a worse
answer context: faithfulness drops by 15.9 percentage points and
hallucination rises from 0\% to 10\%. The mechanism is
over-intervention---the OR gate triggers on 100\% of queries, replacing
adequate passages with narrow fragments that break the connective tissue
between evidence segments. On PubMedQA, HCPC-v1 raises mean retrieval
similarity from 0.653 to 0.678, which would appear as an improvement under
standard evaluation, yet hallucination simultaneously rises from 20.0\% to
26.7\%. This confirms that higher per-chunk similarity can co-occur with
lower answer quality when the intervention fragments context coherence.

The following example from the SQuAD evaluation illustrates the mechanism
concretely.

\begin{quote}
\small
\textbf{Query:} ``What city did Super Bowl 50 take place in?''

\textbf{Baseline} (1024-token chunks, no refinement): The retrieved context
contains a coherent passage covering Super Bowl 50's venue, location, and
surrounding logistics. The generator answers: \emph{``Super Bowl 50 took
place in Santa Clara, California (specifically at Levi's Stadium in the San
Francisco Bay Area).''} Faithfulness: 0.9733.

\textbf{HCPC-v1} (naive refinement): All three chunks are flagged by the OR
gate and split into 256-token sub-chunks. The replacement fragments score
higher on query similarity ($\text{sim} = 0.589$) but the sub-chunk
containing the venue location is discarded in favor of one discussing the
game's branding as ``Super Bowl L.'' The generator responds: \emph{``I
cannot find this information in the provided context. The context only states
that Super Bowl 50 was a game to determine the champion\ldots''}
Faithfulness: 0.4730 (hallucinated).

\textbf{HCPC-v2} (selective refinement): The AND gate does not trigger (both
signals must agree), so the original coherent chunks are preserved.
The generator produces the same correct answer as baseline. Faithfulness:
0.9733.
\end{quote}

\noindent This example shows the paradox in microcosm: refinement improved
the per-chunk similarity score while removing the specific evidence the
generator needed. HCPC-v2 avoided the failure by not intervening at all.

Figure~\ref{fig:paradox} visualizes the disconnect between retrieval
similarity and faithfulness across all conditions.
\input{../figures/paradox_plot}

\FloatBarrier
\subsection{Selective refinement recovers faithfulness through restraint}

The key to HCPC-v2's effectiveness is not a better refinement algorithm but
a more conservative refinement \emph{policy}.
Figure~\ref{fig:retrieval_comparison} compares faithfulness and hallucination
across all major retrieval conditions, and
Figure~\ref{fig:refinement_tradeoff} shows that the faithfulness recovery is
tightly coupled to the reduction in refinement rate.
\input{../figures/retrieval_and_tradeoff}

On SQuAD, HCPC-v2 recovers to near-baseline faithfulness (0.7874 vs.\
0.7995) while eliminating the hallucinations introduced by HCPC-v1. It
achieves this by intervening on only 16.7\% of queries, compared to
HCPC-v1's 100\%. The AND-gated detection and rank protection prevent
refinement on passages that are adequate, preserving the coherence of the
evidence set. On PubMedQA, the gain is smaller (0.5783 to 0.5893) because
the method operates on retrieval sets whose weak coherence originates in the
encoder-corpus mismatch, not in individual weak passages that refinement can
fix.

This behavior validates the central design principle of HCPC-v2: on a strong
retrieval set, the safest refinement strategy is to intervene rarely; on a
weak retrieval set, intervention is sometimes necessary but cannot compensate
fully for poor upstream alignment.

\subsection{Chunk size dominates faithfulness variance}

Variance analysis across the Phase~1--2 ablations reveals an unexpectedly
strong effect of chunk size. On SQuAD, one-way ANOVA yields $F = 19.72$
($p < 0.001$, $\eta^2 = 0.100$) for chunk size, compared to $F = 0.046$
($p = 0.830$, $\eta^2 < 0.001$) for prompt strategy. Among the three
configuration factors tested (chunk size, top-$k$, and prompt strategy),
chunk size accounts for 90.1\% of the between-factor variance
(Appendix~\ref{app:variance_model}). The 256-to-1024 token effect
corresponds to Cohen's $d = 1.41$---a large effect by conventional
standards---reflecting that larger chunks preserve complete sentences,
logical transitions, and redundant evidence that collectively reduce
hallucination risk.

This result has immediate practical implications: it overturns the common
engineering priority of optimizing prompt templates before chunking strategy.
Among the parameters we tested, the faithfulness return from
chunk size selection is approximately 450 times larger than from prompt
style ($\eta^2_{\text{chunk}} / \eta^2_{\text{prompt}} = 0.100 / 0.0002$).

\subsection{Threshold sensitivity reveals a robust operating region}

The Phase~5 threshold sweep confirms that HCPC-v2 does not rely on brittle
parameter choices. Table~\ref{tab:threshold} reports representative SQuAD
operating points across the sweep.

\begin{table}[!htbp]
\centering
\caption{Representative threshold configurations on SQuAD. Faithfulness
remains stable across refinement rates from 7\% to 48\%, degrading only
when refinement becomes aggressive enough to erode context coherence. This
broad plateau means HCPC-v2 does not require precise tuning for deployment.}
\label{tab:threshold}
\small
\begin{tabular}{ccccc}
\toprule
$\tau_{\text{sim}}$ & $\tau_{\text{ce}}$ & Faith. & Halluc. & Refine\% \\
\midrule
0.60 & $-0.05$ & 0.7810 & 3.3\% & 7\%  \\
0.55 & $-0.10$ & 0.7815 & 3.3\% & 17\% \\
0.45 & $-0.20$ & 0.7823 & 3.3\% & 33\% \\
0.40 & $-0.30$ & 0.7809 & 3.3\% & 48\% \\
0.35 & $-0.40$ & 0.7506 & 3.3\% & 68\% \\
\bottomrule
\end{tabular}
\end{table}

Figure~\ref{fig:threshold_plateau} shows a broad performance plateau
followed by a sharp decline once refinement exceeds approximately 50\% of
queries.
\input{../figures/threshold_plateau}
\FloatBarrier

The plateau behavior has a coherence-based explanation. Below 50\%
refinement, the method is modifying passages that are genuinely weak, and
the merge-back step restores contiguity in most cases. Above 50\%, the
method begins refining passages that were adequate, and the cumulative
fragmentation degrades the context faster than merge-back can repair it.

\subsection{Domain mismatch exposes the limits of retrieval-augmented generation}

The contrast between SQuAD and PubMedQA is summarized in
Figure~\ref{fig:dataset_comparison}.

PubMedQA is not simply harder in the conventional sense. Its mean retrieval
similarity is comparable to SQuAD's, yet its faithfulness is consistently
lower and its hallucination rate is higher. The difference lies in context
coherence: the retrieved biomedical abstracts are topically adjacent but not
mutually supportive---they use different terminology for related concepts,
present contradictory evidence from different studies, and lack the narrative
continuity that Wikipedia passages provide. This is exactly the kind of
setting where CCS captures what raw similarity scores miss.

The zero-shot comparison makes the deployment consequence stark.
Table~\ref{tab:zeroshot} shows that RAG helps on SQuAD but actively hurts on
PubMedQA, and Figure~\ref{fig:zeroshot} visualizes the divergence.

\begin{table}[!htbp]
\centering
\caption{Zero-shot versus best RAG configuration. On SQuAD, retrieval
reduces hallucination as expected. On PubMedQA, retrieval with a
general-domain encoder \emph{increases} hallucination relative to answering
without any retrieved context---a fundamental warning for domain-specific
RAG deployment.}
\label{tab:zeroshot}
\small
\begin{tabular}{llcccc}
\toprule
Model & Dataset & ZS Faith. & ZS Halluc. & RAG Faith. & RAG Halluc. \\
\midrule
Mistral & SQuAD    & 0.6895 & 5.0\%  & 0.7775 & 3.0\% \\
Mistral & PubMedQA & 0.6118 & 8.0\%  & 0.6004 & 20.0\% \\
Llama-3 & SQuAD    & 0.6967 & 13.0\% & 0.7752 & 0.0\% \\
Llama-3 & PubMedQA & 0.6110 & 4.0\%  & 0.6233 & 10.0\% \\
\bottomrule
\end{tabular}
\end{table}

\input{../figures/dataset_and_zeroshot}

On PubMedQA, Mistral's hallucination rate rises from 8.0\% (zero-shot) to
20.0\% (best RAG)---a 2.5$\times$ increase. Llama-3's rises from 4.0\% to
10.0\%. This occurs because the general-domain encoder retrieves passages
that are topically related but semantically inconsistent, producing a
low-coherence context that the generator cannot reconcile. The model's
parametric knowledge, while less specific, at least produces internally
consistent answers. This finding argues that coherence should be validated
before any RAG system is deployed in a specialized domain.

\FloatBarrier
